% ============================================================
%  FLARE Reimplementation Report — IEEE Conference Format (IEEEtran)
%  Overleaf: use IEEEtran class (v1.8b+)
% ============================================================

\documentclass[conference]{IEEEtran}

\usepackage[T1]{fontenc}
\usepackage[utf8]{inputenc}
\usepackage{microtype}
\usepackage{cite}
\usepackage{amsmath,amssymb}
\usepackage{graphicx}
\usepackage{booktabs}
\usepackage{array}
\usepackage{multirow}
\usepackage{xcolor}
\usepackage{url}
\usepackage{hyperref}
\usepackage{listings}
\usepackage{balance}
\usepackage{enumitem}

\hypersetup{
  colorlinks=true,
  linkcolor=blue,
  citecolor=blue,
  urlcolor=blue
}

\lstset{
  basicstyle=\footnotesize\ttfamily,
  breaklines=true,
  frame=single,
  backgroundcolor=\color{gray!8},
  keywordstyle=\color{blue}\bfseries,
  commentstyle=\color{green!50!black},
  stringstyle=\color{red!70!black},
  numbers=left,
  numberstyle=\tiny\color{gray},
  numbersep=5pt,
  tabsize=2,
  showstringspaces=false
}

% ============================================================
\begin{document}

\title{%
    Reimplementing FLARE: Active\\
    Retrieval-Augmented Generation\\
    with a Locally Deployed Small Language Model%
}

\author{%
    \IEEEauthorblockN{
        Piyush Satiya -- 202451130
        \hspace{1.5cm}
        Prayag Kirar -- 202451134
        \hspace{1.5cm}
        Sanskar Sharma -- 202451148
    }

    \\

    \IEEEauthorblockA{
        Department of Computer Science and Engineering\\
        Indian Institute of Information Technology Vadodara\\
        Gandhinagar, Gujarat, India
    }

    \vspace{0.3cm}

    \IEEEauthorblockA{
        \textbf{Course Instructor:} Jignesh Patel \quad | \quad \textbf{Course:} CS307
    }
}

\maketitle

% ============================================================
\begin{abstract}
This paper describes our reimplementation of FLARE (Forward-Looking
Active REtrieval Augmented Generation)~\cite{sun2023flare}. FLARE
works by checking how confident the model is about each token it
generates, and only retrieves Wikipedia documents when confidence
drops below a threshold. The original paper used Codex
(${\sim}175$B parameters, closed API). Since Codex is no longer
accessible, we used \textbf{Qwen3-4B-Instruct-2507} with 4-bit NF4
quantization, running fully on a local GPU, paired with a Pyserini
BM25 index over Wikipedia. We tested three modes — No Retrieval,
Single Retrieval, and FLARE — on three benchmark datasets:
2WikiMultiHopQA, StrategyQA, and ASQA. Even with a model roughly
40$\times$ smaller, the main finding holds on all three datasets:
FLARE outperforms both baselines. On StrategyQA the FLARE gain
over Single Retrieval ($+9.2$\% EM) is almost the same as the
paper's ($+8.7$\% EM), which suggests the mechanism works at
smaller scales too. On ASQA the improvement is actually larger
in our case, which we think is because smaller models need
retrieval more.
\end{abstract}

\begin{IEEEkeywords}
Retrieval-Augmented Generation, Active Retrieval, FLARE,
BM25, Qwen, Multi-Hop QA, Open-Domain QA, Large Language Models,
Token Confidence
\end{IEEEkeywords}

% ============================================================
\section{Introduction}
\label{sec:intro}

Language models often generate text that sounds correct but is
factually wrong — a problem commonly called hallucination~\cite{ji2023hallucination}.
Retrieval-Augmented Generation (RAG)~\cite{lewis2020rag} is one
approach to reduce this: before generating an answer, the system
fetches relevant documents from Wikipedia and gives them to the model
as context. This works reasonably well, but the retrieval only happens
once, using the original question as the search query.

The problem is that for multi-step questions, the original question is
not always enough to find the right documents. For example, the question
\emph{``Who is the mother of the director of Polish-Russian War?''} first
requires knowing who directed the film, and only then can you search
for that person's mother. A single retrieval at the beginning cannot
handle this kind of chained reasoning.

FLARE (Sun et al., EMNLP 2023)~\cite{sun2023flare} tries to fix this by
making retrieval happen \emph{during} generation rather than only
at the start. The model generates a tentative next sentence, and if it
is not confident about certain words, those words are masked and the
remaining confident words are used to search Wikipedia. Then the
sentence is regenerated with the retrieved context. This loop
continues until the answer is complete.

We reimplemented this system using open-source tools that run locally,
and checked whether the paper's main claims hold when using a much
smaller model. The code is available at:
\url{https://github.com/PIYUSH-SATIYA/active-rag}.

% ============================================================
\section{Background}
\label{sec:background}

\subsection{Retrieval-Augmented Generation}
The original RAG paper~\cite{lewis2020rag} proposed retrieving
documents from a dense index and using them to condition generation.
FiD (Fusion-in-Decoder)~\cite{izacard2020leveraging} took this
further by fusing multiple retrieved passages in the encoder.
Both methods do retrieval once at query time.

\subsection{Iterative Retrieval}
Some later work tried retrieval at multiple steps.
IRCOT~\cite{trivedi2022interleaving} interleaves chain-of-thought
steps with retrieval calls. Self-Ask~\cite{press2022measuring} has
the model generate explicit sub-questions and retrieves for each.
FLARE is different because it does not ask the model to plan out
sub-questions — instead it watches the token probabilities and
retrieves automatically when uncertainty appears.

\subsection{Using Confidence for Generation}
The idea of using token probabilities as an uncertainty signal is not
entirely new. It has been used in selective
prediction~\cite{kamath2020selective} and for detecting knowledge
conflicts~\cite{longpre2021entity}. What FLARE does differently is
use confidence to \emph{decide when to retrieve}, which is a
practical and fairly elegant application of this idea.

% ============================================================
\section{The FLARE Paper: What It Proposes}
\label{sec:paper_overview}

\subsection{The Generation Loop}
At each step, given question $q$ and what has been generated so far
$s_{1..t-1}$, FLARE does the following:

\begin{enumerate}[leftmargin=*, itemsep=2pt]
  \item Generate the next sentence $\hat{s}_t$ and record the
    probability of each token: $\{p_1, \ldots, p_n\}$.
  \item If all token probabilities are above threshold $\theta$,
    accept the sentence and move on.
  \item Otherwise, drop tokens with $p_i < \beta$ (these might be
    hallucinations) and use the rest as a search query $q_t$.
  \item Retrieve top-$k$ Wikipedia passages $\mathcal{D}_t$ using
    BM25 with query $q_t$.
  \item Regenerate the sentence with $\mathcal{D}_t$ added to the
    context.
  \item Add the sentence to the answer and repeat.
\end{enumerate}

\subsection{Two Variants}
The paper has two versions of FLARE:
\begin{itemize}[leftmargin=*]
  \item \textbf{FLARE\textsubscript{direct}} — uses the look-ahead
    sentence (after masking uncertain tokens) as the search query.
    This is what we implemented.
  \item \textbf{FLARE\textsubscript{instruct}} — explicitly tells
    the model to generate a search query using special prompting.
\end{itemize}
FLARE\textsubscript{direct} gets better numbers in the paper
(51.0\% vs 42.4\% EM on 2WikiMultiHopQA) so we focused on that.

\subsection{Original Setup}
The paper ran everything with Codex (\texttt{text-davinci-003},
${\sim}175$B parameters) through the OpenAI API. They built a custom
BM25 index over Wikipedia, split into 100-word passages. Evaluation
used 500 development examples per dataset with greedy decoding.

% ============================================================
\section{System Architecture}
\label{sec:architecture}

Our system is designed to match the paper's architecture as closely
as possible, given that we had to swap the model and the index.
Figure~\ref{fig:architecture} shows how the parts connect.

\begin{figure}[htbp]
\centering
\fbox{\parbox{0.9\linewidth}{\centering\vspace{1.2cm}
  \small
  \textbf{Question}\\
  $\downarrow$\\
  \textbf{FLAREAgent} (agent.py)\\
  $\updownarrow$\hspace{1.5cm}$\updownarrow$\\
  \textbf{QwenModel}\hspace{0.6cm}\textbf{Retriever}\\
  (qwen.py)\hspace{0.9cm}(retriever.py)\\
  $\updownarrow$\hspace{2.8cm}$\updownarrow$\\
  \textbf{ConfidenceChecker}\quad\textbf{Pyserini BM25}\\
  confidence.py\hspace{0.7cm}\texttt{wikipedia-dpr}\\
  $\downarrow$\\
  \textbf{Final Answer}
  \vspace{1.2cm}
}}
\caption{System architecture. The FLAREAgent runs the loop,
calling the model and retriever as needed each step.}
\label{fig:architecture}
\end{figure}

There are five main components:

\subsubsection*{Language Model (\texttt{src/models/qwen.py})}
This wraps the Qwen3-4B model. The important part is that it
returns per-token probabilities from \texttt{model.generate()},
using the \texttt{output\_scores=True} flag. Without this, the
confidence check cannot work. Output is stored in a
\texttt{GeneratorResult} dataclass with the text, tokens,
and their probabilities.

\subsubsection*{Confidence Checker (\texttt{src/flare/confidence.py})}
Very simple — checks if the minimum token probability in a sentence
is below $\theta$. If yes, retrieval is triggered.

\subsubsection*{Look-ahead Generator (\texttt{src/flare/lookahead.py})}
Builds the prompt and calls the model to get the next sentence.
Uses greedy decoding so results are deterministic.

\subsubsection*{Retriever (\texttt{src/retriever.py})}
Uses Pyserini's \texttt{LuceneSearcher} with the pre-built
\texttt{wikipedia-dpr} index — 21 million 100-word passages.
Same passage size as the paper.

\subsubsection*{FLARE Agent (\texttt{src/flare/agent.py})}
The main class. Implements all three evaluation modes and saves
traces to disk after every step, so we can analyze what happened.

% ============================================================
\section{Implementation Details}
\label{sec:implementation}

\subsection{Model and Quantization}

We could not use Codex since OpenAI removed it. We picked
\textbf{Qwen3-4B-Instruct-2507}~\cite{qwen3} because it is
open-source, reasonably capable, and fits on a single 16\,GB GPU
after quantization. Full FP16 takes about 8\,GB; with 4-bit NF4
quantization via \texttt{bitsandbytes}~\cite{dettmers2024qlora}
it drops to roughly 2.5\,GB:

\begin{lstlisting}[language=Python, caption={NF4 quantization config (src/models/qwen.py)}]
BitsAndBytesConfig(
    load_in_4bit=True,
    bnb_4bit_quant_type="nf4",
    bnb_4bit_compute_dtype=torch.float16,
    bnb_4bit_use_double_quant=True,
)
\end{lstlisting}

Weights are stored in 4-bit but matrix multiplications happen in
FP16, which keeps the quality reasonable.

\subsection{Adapting for a Chat Model}

This was the main design challenge. Codex is a completion model — you
give it text and it continues it. Qwen3-4B-Instruct is a chat model
that expects messages in a specific format. We worked around this by
putting the entire few-shot prompt into a single user message and
running it through Qwen's chat template:

\begin{lstlisting}[language=Python, caption={Wrapping the few-shot prompt as a chat message}]
messages = [{"role": "user",
             "content": full_fewshot_prompt}]
input_ids = tokenizer.apply_chat_template(
    messages, tokenize=True,
    add_generation_prompt=True,
    return_tensors="pt"
)
\end{lstlisting}

This is the one place where our implementation structurally differs
from the paper's prompt setup.

\subsection{Few-Shot Examples}

Like the paper, we use 8 question-answer examples written by hand in
chain-of-thought format. The model sees these before every question:

\begin{lstlisting}[caption={Example from src/templates.py}]
Q: Were Pavel Urysohn and Leonid Levin both Russian?
A: Pavel Urysohn was a Soviet mathematician.
   Leonid Levin is a Soviet-American mathematician.
   So the answer is yes.
\end{lstlisting}

Every answer ends with ``So the answer is X.'' and we extract the
final answer using that pattern. No fine-tuning was done.

\subsection{Wikipedia Retrieval}

We used Pyserini's \texttt{wikipedia-dpr} index~\cite{pyserini}
instead of building our own. It has 21M passages of about 100 words
each, from a 2018 Wikipedia dump — same segmentation as the paper.
Building the index from scratch needs a 20\,GB dump download; the
prebuilt index avoids that.

BM25 scores each passage as:
\begin{equation}
  \text{BM25}(d,q) = \sum_{t \in q}
    \text{IDF}(t) \cdot
    \frac{f(t,d)(k_1+1)}{f(t,d)+k_1(1-b+b\frac{|d|}{\text{avgdl}})}
\end{equation}
with $k_1 = 0.9$ and $b = 0.4$ (Pyserini defaults).

\subsection{Hyperparameters}

We matched the paper's hyperparameters wherever possible.
Table~\ref{tab:hyperparams} lists everything.

\begin{table}[htbp]
\caption{Hyperparameter Configuration}
\label{tab:hyperparams}
\centering
\begin{tabular}{lrl}
\toprule
\textbf{Parameter} & \textbf{Value} & \textbf{Description} \\
\midrule
$\theta$ & 0.8  & Confidence threshold (triggers retrieval) \\
$\beta$  & 0.4  & Masking threshold (query building) \\
$k$      & 2    & Passages retrieved per query \\
Look-ahead tokens & 64 & Max tokens per generation step \\
Max answer tokens & 256 & Total answer length cap \\
Few-shot examples & 8  & In-context demonstrations \\
Max query words   & 64 & Search query length cap \\
Decoding          & Greedy & \texttt{do\_sample=False} \\
\bottomrule
\end{tabular}
\end{table}

\subsection{Sentence Splitting}

FLARE works sentence by sentence. After each look-ahead we cut the
output at the first sentence boundary using a regex:

\begin{lstlisting}[language=Python, caption={Sentence boundary detection}]
_SENT = re.compile(r'(?<=[.!?])\s+')

def _truncate_at_sentence_boundary(text):
    parts = _SENT.split(text.strip(), maxsplit=1)
    return parts[0].strip()
\end{lstlisting}

\subsection{Output Files}

For each experiment the runner saves three files:
\begin{itemize}[leftmargin=*]
  \item \texttt{predictions.jsonl} — question, gold answer, predicted
    answer, EM and F1 per example.
  \item \texttt{metrics.json} — aggregate numbers.
  \item \texttt{traces.jsonl} — every generation step with confidence
    scores, queries sent to BM25, and what came back.
\end{itemize}
If a run crashes, the runner reads back the completed lines and skips
them on restart, so compute is not wasted.

% ============================================================
\section{Datasets}
\label{sec:datasets}

We tested on three of the four datasets from the paper. The fourth
(WikiASP) had a data loading issue that we could not resolve in time
and is not included here. Table~\ref{tab:datasets} gives an overview.

\begin{table}[htbp]
\caption{Datasets Used}
\label{tab:datasets}
\centering
\begin{tabular}{p{2.0cm}p{1.8cm}p{1.5cm}r}
\toprule
\textbf{Dataset} & \textbf{Task Type} & \textbf{Answer Format} & \textbf{Dev Size} \\
\midrule
2WikiMultiHopQA & Multi-hop QA        & Short string  & 500 \\
StrategyQA      & Commonsense yes/no  & true/false    & 500 \\
ASQA            & Ambiguous QA        & Multi-answer  & 500 \\
\bottomrule
\end{tabular}
\end{table}

\subsection{2WikiMultiHopQA}
This dataset~\cite{ho2020constructing} has questions that need
exactly two Wikipedia lookups to answer. For example:
\emph{``Who is the mother of the director of Polish-Russian War?''}
You first need to find the director, then find his mother. The
answers are short strings (names, dates, places), so Exact Match
is a reasonable metric.

\subsection{StrategyQA}
StrategyQA~\cite{geva2021aristotle} is a yes/no dataset where
each question needs implicit multi-step reasoning. For example:
\emph{``Is the tibia necessary to win the Stanley Cup?''} — you
need to know what the tibia is and how the Stanley Cup works.
The gold labels are stored as \texttt{true}/\texttt{false}. Our
model outputs \texttt{yes}/\texttt{no}, so we normalize before
scoring.

\subsection{ASQA}
ASQA~\cite{stelmakh2022asqa} contains ambiguous questions that
have multiple valid answers depending on context or interpretation.
For example: \emph{``Who has the highest goals in world football?''}
— different sources give different answers (Cristiano Ronaldo, Josef
Bican, etc.). Evaluation checks if the model's answer matches any
item in the gold answer set.

% ============================================================
\section{Evaluation Metrics}
\label{sec:metrics}

\subsection{Exact Match}
Both the prediction and gold answer are normalized — lowercase,
punctuation removed, articles (\emph{a, an, the}) dropped,
extra spaces collapsed — and then compared character by character.
A prediction either matches or it doesn't. We report the percentage
of examples that match.

EM is strict. If the model says \emph{``The mother was Małgorzata
Braunek''} instead of just \emph{``Małgorzata Braunek''}, it scores
0, even though the right name is in there.

\subsection{Token-level F1}
Both strings are treated as bags of words and we compute the
harmonic mean of precision and recall:
\begin{equation}
  F_1 = \frac{2 \cdot |P \cap G|}{|P| + |G|}
\end{equation}
where $P$ is the predicted token set and $G$ is the gold token set.
F1 gives partial credit, so it is always at least as high as EM.

\subsection{Metrics We Did Not Implement}
The paper uses Disambig-F1 for ASQA (checks coverage across all
valid interpretations of an ambiguous question) and ROUGE-L plus
UniEval for WikiASP. We only computed EM and F1 due to time and
scope constraints.

% ============================================================
\section{Results}
\label{sec:results}

All experiments use 500 development examples. Paper numbers are from
Sun et al.~\cite{sun2023flare} using Codex.

\subsection{2WikiMultiHopQA}

\begin{table}[htbp]
\caption{Results on 2WikiMultiHopQA}
\label{tab:2wiki}
\centering
\begin{tabular}{lcccc}
\toprule
\multirow{2}{*}{\textbf{Method}} &
  \multicolumn{2}{c}{\textbf{Ours (Qwen3-4B)}} &
  \multicolumn{2}{c}{\textbf{Paper (Codex)}} \\
\cmidrule(lr){2-3}\cmidrule(lr){4-5}
& EM & F1 & EM & F1 \\
\midrule
No Retrieval     & 9.6\%  & 15.9\% & 28.2\% & 36.8\% \\
Single Retrieval & 21.2\% & 25.2\% & 39.4\% & 48.8\% \\
\textbf{FLARE}   & \textbf{23.2\%} & \textbf{28.4\%}
                 & \textbf{51.0\%} & \textbf{59.7\%} \\
\midrule
\multicolumn{5}{l}{\small FLARE gain: Ours $+2.0\%$ EM \quad Paper $+11.6\%$ EM} \\
\bottomrule
\end{tabular}
\end{table}

The ordering No Retrieval $<$ Single Retrieval $<$ FLARE holds in
our results, which matches the paper. The absolute numbers are lower
across the board by roughly 18 percentage points. We think this gap
comes from the model — a 4B quantized model just knows fewer facts
than a 175B model. The FLARE gain here is smaller than the paper's
($+2.0\%$ vs $+11.6\%$ EM). This is partly because our model
retrieves at nearly every step (88.7\% of steps), which introduces
noise from irrelevant passages and reduces the net benefit of
retrieval. This is discussed more in Section~\ref{sec:analysis}.

\subsection{StrategyQA}

\begin{table}[htbp]
\caption{Results on StrategyQA}
\label{tab:stratqa}
\centering
\begin{tabular}{lccc}
\toprule
\textbf{Method} & \textbf{Ours EM} & \textbf{Paper EM} & \textbf{Gap} \\
\midrule
No Retrieval     & 15.6\% & 72.9\% & $-57.3$ pp \\
Single Retrieval & 31.0\% & 68.6\% & $-37.6$ pp \\
\textbf{FLARE}   & \textbf{40.2\%} & \textbf{77.3\%} & $-37.1$ pp \\
\midrule
\multicolumn{4}{l}{\small FLARE gain: Ours $+9.2\%$ \quad Paper $+8.7\%$ (very close)} \\
\bottomrule
\end{tabular}
\end{table}

This is probably our best result. The FLARE improvement over Single
Retrieval is $+9.2\%$ EM in our case and $+8.7\%$ EM in the paper.
That is close enough that we feel the mechanism is working correctly
at this scale.

The overall numbers are much lower than the paper (40.2\% vs 77.3\%
for FLARE), which is expected given the model size difference. These
questions need commonsense reasoning that smaller models tend to
struggle with.

One thing worth noting: in the paper, Single Retrieval is actually
slightly worse than No Retrieval on this dataset (68.6\% vs 72.9\%).
We do not see that — in our case Single Retrieval helps. This likely
happens because our model has less built-in knowledge and benefits
from any grounding, whereas Codex already knows most of these
answers and extra retrieved documents can sometimes confuse it.

\subsection{ASQA}

\begin{table}[htbp]
\caption{Results on ASQA}
\label{tab:asqa}
\centering
\begin{tabular}{lcccc}
\toprule
\multirow{2}{*}{\textbf{Method}} &
  \multicolumn{2}{c}{\textbf{Ours (Qwen3-4B)}} &
  \multicolumn{2}{c}{\textbf{Paper (Codex)}} \\
\cmidrule(lr){2-3}\cmidrule(lr){4-5}
& EM & F1 & EM & Disambig-F1 \\
\midrule
No Retrieval     & 15.0\% & 22.2\% & 33.8\% & 24.2\% \\
Single Retrieval & 15.0\% & 24.9\% & 40.0\% & 27.1\% \\
\textbf{FLARE}   & \textbf{29.0\%} & \textbf{37.6\%}
                 & \textbf{41.3\%} & \textbf{28.2\%} \\
\midrule
\multicolumn{5}{l}{\small FLARE gain: Ours $+14.0\%$ EM \quad Paper $+1.3\%$ EM} \\
\bottomrule
\end{tabular}
\end{table}

The ASQA result is interesting. Our FLARE improvement ($+14.0\%$ EM)
is much larger than the paper's ($+1.3\%$ EM). Single Retrieval
does not help at all over No Retrieval in our case (both 15.0\% EM)
— which makes sense because fetching documents for an ambiguous
question just gives you documents about one interpretation. FLARE
can refine what it is looking for over multiple steps, which is
exactly what these questions need.

We think this larger gain happens because our model has less
internal knowledge than Codex, so the iterative retrieval fills a
bigger gap. Codex can often just answer from memory; Qwen3-4B needs
the retrieval to compensate.

The F1 scores also show a clear pattern even when EM does not
differentiate as much. Single Retrieval's F1 (24.9\%) is higher
than No Retrieval (22.2\%) even though EM is the same, which
suggests the retrieved documents do help with partial correctness.

\subsection{FLARE Gain Summary}

Table~\ref{tab:gain_summary} shows FLARE's improvement over Single
Retrieval across all three datasets. This is the most direct way
to evaluate whether the FLARE mechanism itself is adding value.

\begin{table}[htbp]
\caption{FLARE Gain over Single Retrieval (EM)}
\label{tab:gain_summary}
\centering
\begin{tabular}{p{2.2cm}ccc}
\toprule
\textbf{Dataset} & \textbf{Ours} & \textbf{Paper} & \textbf{Assessment} \\
\midrule
2WikiMultiHopQA & $+2.0$\%  & $+11.6$\% & Direction reproduced \\
StrategyQA      & $+9.2$\%  & $+8.7$\%  & \checkmark~Very close \\
ASQA            & $+14.0$\% & $+1.3$\%  & Direction reproduced \\
\bottomrule
\end{tabular}
\end{table}

% ============================================================
\section{FLARE Behavior Analysis}
\label{sec:analysis}

We looked at the \texttt{traces.jsonl} from the 2WikiMultiHopQA
run to understand what the FLARE loop is actually doing step by step.
Table~\ref{tab:traces} shows the numbers.

\begin{table}[htbp]
\caption{FLARE Loop Statistics -- 2WikiMultiHopQA (500 examples)}
\label{tab:traces}
\centering
\begin{tabular}{lr}
\toprule
\textbf{Metric}                          & \textbf{Value}  \\
\midrule
Total generation steps                   & 2{,}242         \\
Average steps per example                & 4.48            \\
Steps where model was confident (no retrieval) & 253\ (11.3\%) \\
Steps where retrieval was triggered      & 1{,}989\ (88.7\%) \\
Average retrievals per example           & 3.98            \\
Mean min token probability per step      & 0.417           \\
Median min token probability             & 0.391           \\
\bottomrule
\end{tabular}
\end{table}

The retrieval rate of 88.7\% is the main issue. FLARE is supposed to
retrieve \emph{selectively} — the paper says 40--80\% is the useful
range. At 88.7\% we are retrieving at almost every step, which means
we are sometimes injecting Wikipedia passages that are not relevant,
adding noise to the generation.

The reason this happens is that our model produces much lower token
probabilities than Codex, even when it knows the answer. The mean
minimum probability per step is 0.417, which is well below the
threshold of 0.8. A single token at 0.79 triggers retrieval, and
for Qwen3-4B that happens almost always.

The fix would be to lower $\theta$ to around 0.5 for this model
size, which would bring the retrieval rate into the optimal range.
We did not tune this in the current experiments, but it is probably
why the 2WikiMH gain is smaller than expected. FLARE still beats
both baselines despite the over-retrieval, which suggests the
algorithm is fundamentally working.

% ============================================================
\section{Differences from the Paper}
\label{sec:differences}

Table~\ref{tab:differences} lists the main ways our implementation
differs from Sun et al.

\begin{table}[htbp]
\caption{Implementation Differences from Sun et al.~\cite{sun2023flare}}
\label{tab:differences}
\centering
\begin{tabular}{p{1.5cm}p{2.2cm}p{2.2cm}}
\toprule
\textbf{Aspect}       & \textbf{Paper}              & \textbf{Ours}               \\
\midrule
LM Model      & Codex $\sim$175B (API)      & Qwen3-4B (local, 4-bit)     \\
LM Type       & Completion                  & Chat / Instruct             \\
Precision     & FP16                        & 4-bit NF4                   \\
Wiki index    & Custom BM25 dump            & Pyserini wikipedia-dpr      \\
FLARE variant & direct + instruct           & direct only                 \\
ASQA metric   & Disambig-F1 + ROUGE-L       & EM + F1                     \\
Datasets      & 4 (incl. WikiASP)           & 3 (WikiASP excluded)        \\
\bottomrule
\end{tabular}
\end{table}

The biggest difference by far is the model. Everything else —
the algorithm, the threshold values, the passage size, the
BM25 retrieval method — is kept as close to the paper as possible.
The chat-template adaptation for Qwen adds some formatting overhead
but does not change how the retrieval logic works.

% ============================================================
\section{Discussion}
\label{sec:discussion}

\subsection{The Main Finding Holds}
Across all three datasets we tested, FLARE consistently beats both
No Retrieval and Single Retrieval. This is the central claim of the
paper and it holds even with a model that is around 40$\times$ smaller.
That is a meaningful result.

\subsection{Smaller Models Can Benefit More}
The ASQA result is somewhat unexpected. We get a $+14\%$ EM improvement
from FLARE whereas the paper only gets $+1.3\%$. One interpretation
is that smaller models benefit more from active retrieval because they
cannot rely on stored knowledge — they \emph{need} the Wikipedia
passages to fill in gaps. Codex can often answer from memory, but
Qwen3-4B cannot, so the retrieval loop helps it more.

\subsection{Threshold Needs Re-tuning}
Using $\theta = 0.8$ for a 4B model is probably not the right choice.
The paper tuned this value on Codex, and Codex has a very different
probability distribution. A threshold around 0.5 would likely work
better here and should improve the 2WikiMH numbers specifically. This
is the most obvious next step.

\subsection{Limitations}
\begin{itemize}[leftmargin=*]
  \item WikiASP was not evaluated due to a data loading issue with
    the HuggingFace parquet schema.
  \item We did not implement FLARE\textsubscript{instruct}.
  \item Only the development split was used, not test.
  \item We did not implement Disambig-F1 or ROUGE-L metrics.
  \item All results are from a single run; no statistical testing
    was done.
\end{itemize}

% ============================================================
\section{Conclusion}
\label{sec:conclusion}

We reimplemented FLARE using a local 4-bit quantized model
(Qwen3-4B-Instruct-2507) and a pre-built Pyserini BM25 index over
Wikipedia. Even with a model roughly 40$\times$ smaller than the
original, the core pattern -- FLARE outperforming both No Retrieval
and Single Retrieval -- holds on all three datasets we evaluated.

On StrategyQA the FLARE gain ($+9.2\%$ EM) nearly matches the paper's
($+8.7\%$ EM), showing the mechanism is consistent across model
scales. On ASQA our improvement is actually larger than the paper's
($+14.0\%$ vs $+1.3\%$ EM), which suggests smaller models benefit
more from active retrieval when they lack the parametric knowledge to
answer without help.

The main open issue is threshold calibration. The $\theta = 0.8$
setting causes over-retrieval (88.7\% rate) with this smaller model,
which reduces the benefit on multi-hop questions. Lowering the
threshold to around 0.5 would likely improve results without changing
anything in the algorithm.

Overall, the key ideas from the FLARE paper can be reproduced using
open-source tools that run on consumer hardware, without needing
access to large proprietary APIs.

% ============================================================
\section*{Acknowledgments}

The authors thank \textbf{Jignesh Patel sir} for guidance and feedback
throughout this project, and the open-source communities behind
HuggingFace Transformers, Pyserini, and the Qwen model family.

% ============================================================
\begin{thebibliography}{99}

\bibitem{sun2023flare}
Z.~Sun, X.~Wang, Y.~Tay, Y.~Yang, and D.~Zhou,
``Active retrieval augmented generation,''
in \emph{Proc.\ Conf.\ Empirical Methods in Natural Language Processing
  (EMNLP)}, Singapore, Dec.\ 2023, pp.\ 7969--7992.

\bibitem{lewis2020rag}
P.~Lewis, E.~Perez, A.~Piktus, F.~Petroni, V.~Karpukhin,
N.~Goyal, H.~K\"{u}ttler, M.~Lewis, W.-T.~Yih, T.~Rockt\"{a}schel,
S.~Riedel, and D.~Kiela,
``Retrieval-augmented generation for knowledge-intensive NLP tasks,''
in \emph{Advances in Neural Information Processing Systems (NeurIPS)},
vol.~33, pp.\ 9459--9474, 2020.

\bibitem{izacard2020leveraging}
G.~Izacard and E.~Grave,
``Leveraging passage retrieval with generative models for open domain
  question answering,''
in \emph{Proc.\ Conf.\ European Chapter of the ACL (EACL)},
2021, pp.\ 874--880.

\bibitem{trivedi2022interleaving}
H.~Trivedi, N.~Balasubramanian, T.~Khot, and A.~Sabharwal,
``Interleaving retrieval with chain-of-thought reasoning for
  knowledge-intensive multi-step questions,''
in \emph{Proc.\ Annual Meeting of the ACL}, Toronto, 2023,
pp.\ 10014--10037.

\bibitem{press2022measuring}
O.~Press, M.~Zhang, S.~Min, L.~Schmidt, N.~A.~Smith, and M.~Lewis,
``Measuring and narrowing the compositionality gap in language models,''
in \emph{Findings of EMNLP}, 2023.

\bibitem{ji2023hallucination}
Z.~Ji, N.~Lee, R.~Frieske, T.~Yu, D.~Su, Y.~Xu,
E.~Ishii, Y.~J.~Bang, A.~Madotto, and P.~Fung,
``Survey of hallucination in natural language generation,''
\emph{ACM Computing Surveys}, vol.~55, no.~12, pp.\ 1--38, 2023.

\bibitem{kamath2020selective}
A.~Kamath, R.~Jia, and P.~Liang,
``Selective question answering under domain shift,''
in \emph{Proc.\ Annual Meeting of the ACL}, 2020, pp.\ 5684--5696.

\bibitem{longpre2021entity}
S.~Longpre, K.~Perisetla, A.~Chen, N.~Ramesh,
C.~DuBois, and S.~Singh,
``Entity-based knowledge conflicts in question answering,''
in \emph{Proc.\ EMNLP}, 2021, pp.\ 7052--7063.

\bibitem{ho2020constructing}
X.~Ho, A.-K.~D.~Nguyen, S.~Sugawara, and A.~Aizawa,
``Constructing a multi-hop QA dataset for comprehensive evaluation of
  reasoning steps,''
in \emph{Proc.\ COLING}, 2020, pp.\ 6609--6625.

\bibitem{geva2021aristotle}
M.~Geva, D.~Khashabi, E.~Segal, T.~Khot, D.~Roth, and J.~Berant,
``Did Aristotle use a laptop? A question answering benchmark with
  implicit reasoning strategies,''
\emph{Transactions of the ACL}, vol.~9, pp.\ 346--361, 2021.

\bibitem{stelmakh2022asqa}
I.~Stelmakh, Y.~Luan, B.~Dhingra, and M.-W.~Chang,
``ASQA: Factoid questions meet long-form answers,''
in \emph{Proc.\ EMNLP}, 2022, pp.\ 8273--8288.

\bibitem{hayashi2021wikiasp}
H.~Hayashi, P.~Budania, P.~Wang, C.~Ackerson, R.~Neervannan, and
G.~Neubig,
``WikiAsp: A dataset for multi-domain aspect-based summarization,''
\emph{Transactions of the ACL}, vol.~9, pp.\ 211--225, 2021.

\bibitem{pyserini}
J.~Lin, X.~Ma, S.-C.~Lin, J.~Yang, R.~Pradeep, and R.~Nogueira,
``Pyserini: A Python toolkit for reproducible information retrieval
  research with sparse and dense representations,''
in \emph{Proc.\ ACM SIGIR}, 2021, pp.\ 2356--2362.

\bibitem{qwen3}
Qwen Team, Alibaba Cloud,
``Qwen3 Technical Report,''
\emph{arXiv preprint arXiv:2505.09388}, 2025.
[Online]. Available: \url{https://arxiv.org/abs/2505.09388}

\bibitem{dettmers2024qlora}
T.~Dettmers, A.~Pagnoni, A.~Holtzman, and L.~Zettlemoyer,
``QLoRA: Efficient finetuning of quantized LLMs,''
in \emph{Advances in Neural Information Processing Systems (NeurIPS)},
vol.~36, 2024.

\end{thebibliography}

\balance
\end{document}
